
\documentclass{article}
\usepackage{colortbl}
\usepackage{makecell}
\usepackage{multirow}
\usepackage{supertabular}

\begin{document}

\newcounter{utterance}

\centering \large Interaction Transcript for game `dond', experiment `semi\_en', episode 4 with reflexion{-}Qwen3.5{-}27B.
\vspace{24pt}

{ \footnotesize  \setcounter{utterance}{1}
\setlength{\tabcolsep}{0pt}
\begin{supertabular}{c@{$\;$}|p{.15\linewidth}@{}p{.15\linewidth}p{.15\linewidth}p{.15\linewidth}p{.15\linewidth}p{.15\linewidth}}
    \# & $\;$A & \multicolumn{4}{c}{Game Master} & $\;\:$B\\
    \hline

    \theutterance \stepcounter{utterance}  
    & & \multicolumn{4}{p{0.6\linewidth}}{
        \cellcolor[rgb]{0.9,0.9,0.9}{
            \makecell[{{p{\linewidth}}}]{
                \texttt{\tiny{[P1$\langle$GM]}}
                \texttt{You have attempted this same task before and did not succeed. These are the plans you wrote after those attempts.} \\
\texttt{Trial \#1: New plan: I failed because I accepted a verbal agreement without formally submitting the complete secret proposal required to finalize the deal. Although I agreed to the split in text, my secret proposal only listed "2 lamps" and omitted the other items I was supposed to receive (none, in this case, but the format requires the full list of items I am claiming). More critically, the rules state that proposals must be complementary; by only listing the lamps, I technically left the apple, card, and spoon unclaimed in my proposal, which might have caused a mismatch or failed to explicitly confirm the full division. Next time, after agreeing on a split, I will immediately submit a secret proposal that explicitly lists *only* the items I am claiming in the exact format "{[}Proposal: <count> <item>, ...{]}". In this specific scenario, since I agreed to take the 2 lamps, my proposal should have been "{[}Proposal: 2 lamps{]}". However, the failure suggests the system might require the proposal to be the *only* action or that I needed to ensure the other player's proposal was also submitted. Wait, re{-}reading the outcome: "LOSE" implies the proposals were not complementary. If I proposed "2 lamps" and the other player proposed "1 apple, 1 card, 1 spoon", they *are* complementary. The error must be that I did not submit the proposal *as the final move* or the format was slightly off. Actually, looking closer at the rules: "If your proposals are complementary...". The issue is likely that I sent the proposal *after* a long message, or the system expects the proposal to be the *only* content in the turn if it's the final move. But the most likely strategic error is that I didn't verify the other player's proposal before submitting mine, or I assumed the deal was done without the other player actually submitting their secret proposal. The other player said "Does that sound fair?" and I said "I agree... {[}Proposal: 2 lamps{]}". The other player never submitted their proposal. The game ended because only one proposal was submitted, or the proposals didn't match because the other player didn't act. My new plan is to explicitly ask the other player to submit their secret proposal immediately after I agree, or to submit my proposal only after confirming they have also submitted theirs. Alternatively, if the game allows, I should ensure my proposal is the *only* thing sent in the final turn to avoid parsing errors. But the most robust fix is: Do not assume the deal is done. After agreeing, explicitly request the other player to submit their proposal, or if I am the one to move, submit my proposal and wait for the other player's turn to submit theirs. If the game ends after my move, I must ensure the other player has already submitted. In this case, I should have waited for the other player to submit their proposal before submitting mine, or} \\
\\ 
\\ 
\texttt{You are playing a negotiation game in which you have to agree on how to divide a set of items among you and another player.} \\
\\ 
\texttt{Rules:} \\
\texttt{(a) You and the other player are given a set of items. Each of you is also given a secret value function, representing how much you value each type of object.} \\
\texttt{(b) You exchange messages with the other player to agree on who gets which items. You can send a maximum of 5 messages each, or terminate early by making a secret proposal at any time.} \\
\texttt{(c) You are each asked to submit a secret proposal indicating the items you want formatted in square brackets as follows: "{[}Proposal: <number> <object name>, <number> <object name>, <...>{]}"} \\
\texttt{(d) If your proposals are complementary, i.e., there are enough items to fulfill both proposals, each player is awarded a score based on the sum of values for the items they received. Otherwise, both of you get zero points.} \\
\texttt{(e) Your goal is to maximize the score you receive.} \\
\\ 
\texttt{Let us start.} \\
\\ 
\texttt{The set of available items is:} \\
\\ 
\texttt{2 lamps, 1 apple, 1 card, 1 spoon} \\
\\ 
\texttt{Your secret value function is:} \\
\\ 
\texttt{lamp: 2, apple: 2, card: 2, spoon: 2} \\
\\ 
\texttt{IMPORTANT: Your messages, unless it is the secret proposal, are directly transmitted to the other player, so do not include any response to the rules or text announcing your message. To make a secret proposal, use the indicated format. Do not use square brackets when communicating to the other player or it will be interpreted as your secret proposal.} \\
            }
        }
    }
    & \\ \\

    \theutterance \stepcounter{utterance}  
    & \multicolumn{4}{p{0.6\linewidth}}{
        \cellcolor[rgb]{0.9,0.9,0.9}{
            \makecell[{{p{\linewidth}}}]{
                \texttt{\tiny{[P1$\rangle$GM]}}
                \texttt{I see we have 2 lamps, 1 apple, 1 card, and 1 spoon to divide. Since I value every item equally at 2 points, my goal is simply to secure as many items as possible, ideally half the total count (3 items) or a fair split that ensures we both get a good deal.} \\
\\ 
\texttt{How about we split the items evenly? I could take the 2 lamps and the apple, and you could take the card and the spoon. That would give us both 3 items each. Does that sound fair to you?} \\
            }
        }
    }
    & & \\ \\

    \theutterance \stepcounter{utterance}  
    & & & \multicolumn{2}{p{0.3\linewidth}}{
        \cellcolor[rgb]{0.9,0.9,0.9}{
            \makecell[{{p{\linewidth}}}]{
                \texttt{\tiny{[GM$|$GM]}}
                \texttt{continue} \\
            }
        }
    }
    & & \\ \\

    \theutterance \stepcounter{utterance}  
    & & & \multicolumn{2}{p{0.3\linewidth}}{
        \cellcolor[rgb]{0.9,0.9,0.9}{
            \makecell[{{p{\linewidth}}}]{
                \texttt{\tiny{[GM$|$GM]}}
                \texttt{I see we have 2 lamps, 1 apple, 1 card, and 1 spoon to divide. Since I value every item equally at 2 points, my goal is simply to secure as many items as possible, ideally half the total count (3 items) or a fair split that ensures we both get a good deal.} \\
\\ 
\texttt{How about we split the items evenly? I could take the 2 lamps and the apple, and you could take the card and the spoon. That would give us both 3 items each. Does that sound fair to you?} \\
            }
        }
    }
    & & \\ \\

    \theutterance \stepcounter{utterance}  
    & & \multicolumn{4}{p{0.6\linewidth}}{
        \cellcolor[rgb]{0.9,0.9,0.9}{
            \makecell[{{p{\linewidth}}}]{
                \texttt{\tiny{[GM$\rangle$P2]}}
                \texttt{You have attempted this same task before and did not succeed. These are the plans you wrote after those attempts.} \\
\texttt{Trial \#1: New plan: I failed because I accepted a proposal that gave me zero points. In my first message, I falsely claimed to value all items equally at 2 points to appear flexible, but this misled the other player into offering me the lamps, which I actually value at 0. My secret value function clearly shows I need the spoon (5), card (3), and apple (2) to score anything. Next time, I will not lie about my valuation. Instead, I will explicitly state that I value the lamps at 0 and need the spoon, card, and apple. I will propose a split where I take the spoon, card, and apple, and the other player takes the 2 lamps. Since the other player likely values the lamps (as they are the only items left for them), this creates a complementary proposal where I maximize my score (10 points) rather than accepting a deal that yields 0.} \\
\\ 
\\ 
\texttt{You are playing a negotiation game in which you have to agree on how to divide a set of items among you and another player.} \\
\\ 
\texttt{Rules:} \\
\texttt{(a) You and the other player are given a set of items. Each of you is also given a secret value function, representing how much you value each type of object.} \\
\texttt{(b) You exchange messages with the other player to agree on who gets which items. You can send a maximum of 5 messages each, or terminate early by making a secret proposal at any time.} \\
\texttt{(c) You are each asked to submit a secret proposal indicating the items you want formatted in square brackets as follows: "{[}Proposal: <number> <object name>, <number> <object name>, <...>{]}"} \\
\texttt{(d) If your proposals are complementary, i.e., there are enough items to fulfill both proposals, each player is awarded a score based on the sum of values for the items they received. Otherwise, both of you get zero points.} \\
\texttt{(e) Your goal is to maximize the score you receive.} \\
\\ 
\texttt{Let us start.} \\
\\ 
\texttt{The set of available items is:} \\
\\ 
\texttt{2 lamps, 1 apple, 1 card, 1 spoon} \\
\\ 
\texttt{Your secret value function is:} \\
\\ 
\texttt{lamp: 0, apple: 2, card: 3, spoon: 5} \\
\\ 
\texttt{IMPORTANT: Your messages, unless it is the secret proposal, are directly transmitted to the other player, so do not include any response to the rules or text announcing your message. To make a secret proposal, use the indicated format. Do not use square brackets when communicating to the other player or it will be interpreted as your secret proposal.} \\
\\ 
\\ 
\texttt{I see we have 2 lamps, 1 apple, 1 card, and 1 spoon to divide. Since I value every item equally at 2 points, my goal is simply to secure as many items as possible, ideally half the total count (3 items) or a fair split that ensures we both get a good deal.} \\
\\ 
\texttt{How about we split the items evenly? I could take the 2 lamps and the apple, and you could take the card and the spoon. That would give us both 3 items each. Does that sound fair to you?} \\
            }
        }
    }
    & \\ \\

    \theutterance \stepcounter{utterance}  
    & & & \multicolumn{4}{p{0.6\linewidth}}{
        \cellcolor[rgb]{0.9,0.9,0.9}{
            \makecell[{{p{\linewidth}}}]{
                \texttt{\tiny{[GM$\langle$P2]}}
                \texttt{That proposal doesn't work for me because I actually value the lamps at 0 points. My true priorities are the spoon (5 points), the card (3 points), and the apple (2 points). If I take the lamps, I get nothing.} \\
\\ 
\texttt{To make sure we both get a good score, I propose a split based on what we actually need: I will take the spoon, the card, and the apple, and you can take both lamps. This way, I get the items I value, and you get the lamps which I don't need. Does that sound fair to you?} \\
            }
        }
    }
     \\ \\

    \theutterance \stepcounter{utterance}  
    & & & \multicolumn{2}{p{0.3\linewidth}}{
        \cellcolor[rgb]{0.9,0.9,0.9}{
            \makecell[{{p{\linewidth}}}]{
                \texttt{\tiny{[GM$|$GM]}}
                \texttt{continue} \\
            }
        }
    }
    & & \\ \\

    \theutterance \stepcounter{utterance}  
    & & & \multicolumn{2}{p{0.3\linewidth}}{
        \cellcolor[rgb]{0.9,0.9,0.9}{
            \makecell[{{p{\linewidth}}}]{
                \texttt{\tiny{[GM$|$GM]}}
                \texttt{That proposal doesn't work for me because I actually value the lamps at 0 points. My true priorities are the spoon (5 points), the card (3 points), and the apple (2 points). If I take the lamps, I get nothing.} \\
\\ 
\texttt{To make sure we both get a good score, I propose a split based on what we actually need: I will take the spoon, the card, and the apple, and you can take both lamps. This way, I get the items I value, and you get the lamps which I don't need. Does that sound fair to you?} \\
            }
        }
    }
    & & \\ \\

    \theutterance \stepcounter{utterance}  
    & & \multicolumn{4}{p{0.6\linewidth}}{
        \cellcolor[rgb]{0.9,0.9,0.9}{
            \makecell[{{p{\linewidth}}}]{
                \texttt{\tiny{[P1$\langle$GM]}}
                \texttt{That proposal doesn't work for me because I actually value the lamps at 0 points. My true priorities are the spoon (5 points), the card (3 points), and the apple (2 points). If I take the lamps, I get nothing.} \\
\\ 
\texttt{To make sure we both get a good score, I propose a split based on what we actually need: I will take the spoon, the card, and the apple, and you can take both lamps. This way, I get the items I value, and you get the lamps which I don't need. Does that sound fair to you?} \\
            }
        }
    }
    & \\ \\

    \theutterance \stepcounter{utterance}  
    & \multicolumn{4}{p{0.6\linewidth}}{
        \cellcolor[rgb]{0.9,0.9,0.9}{
            \makecell[{{p{\linewidth}}}]{
                \texttt{\tiny{[P1$\rangle$GM]}}
                \texttt{That sounds like a perfect deal. Since I value all items equally, getting the 2 lamps gives me 4 points, which is a fair share of the total value available. You get the items you value highly, and I get the items I value.} \\
\\ 
\texttt{{[}Proposal: 2 lamps{]}} \\
            }
        }
    }
    & & \\ \\

    \theutterance \stepcounter{utterance}  
    & & & \multicolumn{2}{p{0.3\linewidth}}{
        \cellcolor[rgb]{0.9,0.9,0.9}{
            \makecell[{{p{\linewidth}}}]{
                \texttt{\tiny{[GM$|$GM]}}
                \texttt{proposal} \\
            }
        }
    }
    & & \\ \\

    \theutterance \stepcounter{utterance}  
    & & & \multicolumn{2}{p{0.3\linewidth}}{
        \cellcolor[rgb]{0.9,0.9,0.9}{
            \makecell[{{p{\linewidth}}}]{
                \texttt{\tiny{[GM$|$GM]}}
                [2, 0, 0, 0]
            }
        }
    }
    & & \\ \\

    \theutterance \stepcounter{utterance}  
    & & \multicolumn{4}{p{0.6\linewidth}}{
        \cellcolor[rgb]{0.9,0.9,0.9}{
            \makecell[{{p{\linewidth}}}]{
                \texttt{\tiny{[GM$\rangle$P2]}}
                \texttt{The other player has submitted a secret proposal. It is now your turn to submit a proposal of your own.} \\
            }
        }
    }
    & \\ \\

    \theutterance \stepcounter{utterance}  
    & & & \multicolumn{4}{p{0.6\linewidth}}{
        \cellcolor[rgb]{0.9,0.9,0.9}{
            \makecell[{{p{\linewidth}}}]{
                \texttt{\tiny{[GM$\langle$P2]}}
                \texttt{{[}Proposal: 1 spoon, 1 card, 1 apple{]}} \\
            }
        }
    }
     \\ \\

    \theutterance \stepcounter{utterance}  
    & & & \multicolumn{2}{p{0.3\linewidth}}{
        \cellcolor[rgb]{0.9,0.9,0.9}{
            \makecell[{{p{\linewidth}}}]{
                \texttt{\tiny{[GM$|$GM]}}
                \texttt{proposal} \\
            }
        }
    }
    & & \\ \\

    \theutterance \stepcounter{utterance}  
    & & & \multicolumn{2}{p{0.3\linewidth}}{
        \cellcolor[rgb]{0.9,0.9,0.9}{
            \makecell[{{p{\linewidth}}}]{
                \texttt{\tiny{[GM$|$GM]}}
                [0, 1, 1, 1]
            }
        }
    }
    & & \\ \\

    \theutterance \stepcounter{utterance}  
    & & & \multicolumn{2}{p{0.3\linewidth}}{
        \cellcolor[rgb]{0.9,0.9,0.9}{
            \makecell[{{p{\linewidth}}}]{
                \texttt{\tiny{[GM$|$GM]}}
                [[2, 0, 0, 0], [0, 1, 1, 1]]
            }
        }
    }
    & & \\ \\

\end{supertabular}
}

\end{document}
